
\chapter{Supplementary Prompts and Case Studies for Online Reasoning Improvement}
\label{app:p4_supplementary}

This appendix presents detailed execution case studies
illustrating the evidence-grounded reasoning mechanisms
introduced in Part A of the proposed framework.

Three framework variants are examined:
evidence-grounded Reflection (V1),
Reflection with ExpeL-based experiential learning (V2),
and Reflection with ExpeL and the Runtime Evidence
Graph (V3).

Each case study uses recorded experimental outputs
to illustrate how architectural predictions are
extracted, evaluated against simulation evidence,
and incorporated into subsequent reasoning.

The examples are intended to explain the operation
and limitations of the individual mechanisms rather
than establish their aggregate effectiveness,
which is evaluated separately in the main
Experimental Results section.

\section{LLM Prompt Templates}
\label{app:p4_prompts}


\subsection{Evidence-Grounded Reflection Prompt (V1)}
\label{app:p4_reflection_prompt}

V1 introduces a Reflexion-inspired mechanism that
examines LLM-generated architectural reasoning against
empirical runtime evidence collected during DSE.

The Reflection mechanism is implemented as an LLM-based
operation. Its purpose is not to generate new architectural
knowledge, but to identify weaknesses in previously
generated reasoning and produce corrective feedback
for subsequent iterations.

The prompt consists of two components: a method
definition describing the role of Reflexion and
task-specific instructions defining how architectural
claims and their associated runtime evidence must
be evaluated.

The prompt explicitly prioritizes measured simulation
results over previous LLM-generated predictions and
requires the model to distinguish parameter-isolated
observations from comparisons involving simultaneous
architectural changes.

A single Reflection invocation generates two
complementary outputs: claim-level feedback for
individual architectural predictions and an
iteration-level assessment summarizing reasoning
strengths, weaknesses, and corrective guidance.

Listing~\ref{lst:p4_reflection_prompt} presents the
Reflection prompt template used in V1. The fixed
instructions are reproduced from the implementation,
with line wrapping adjusted for readability. The
iteration-specific runtime input is represented by
a placeholder.

\begin{lstlisting}[
    style=remodeljson,
    caption={Reflection prompt template used for
    evidence-grounded reasoning feedback in V1.},
    label={lst:p4_reflection_prompt},
    breaklines=true,
    basicstyle=\ttfamily\footnotesize
]
REFLEXION-INSPIRED REASONING FEEDBACK

In this framework, Reflexion means:

1. Examine the LLM's original extracted reasoning claims.
2. Compare those claims with the runtime evidence accumulated from executed
   designs and measured simulation results.
3. Identify which parts of the reasoning are:
   - directly supported by runtime evidence,
   - contradicted by runtime evidence,
   - unchanged in the observed comparison,
   - non-verifiable,
   - or currently unsupported because evidence is insufficient.
4. Generate verbal corrective feedback that helps the LLM reason more
   accurately in the next iteration.
5. Preserve uncertainty and clearly distinguish:
   - measured observations,
   - evidence-based interpretations,
   - and unsupported hypotheses.

Reflexion does not create new hardware knowledge.
Reflexion does not use the offline benchmark database.
Reflexion does not modify the original extracted claims.
Its output is corrective feedback for future reasoning.


You are now performing the Reflexion-inspired learning step.

The goal of this step is NOT to generate new hardware knowledge.

The goal is to examine the LLM's previous reasoning against empirical
runtime evidence collected during execution and produce corrective
feedback that improves the reasoning of the next iteration.

You must analyze the original extracted claims against the supplied
runtime evidence.

Evidence authority:
1. Measured simulation results are the highest-authority source.
2. Exact parameter-isolated comparisons are stronger than comparisons
   where multiple parameters changed.
3. Previous LLM claims, reflections, and lessons are not ground truth.
4. No offline database, optimal design, design closeness, or external
   architectural knowledge may be used.

For every claim:

- If supporting_evidence exists, explain exactly what was observed.
- If contradicting_evidence exists, explain the contradiction.
- If unchanged_evidence exists, state that no measurable effect was observed.
- If other parameters changed, explicitly state that causality cannot
  be assigned solely to the claimed parameter.
- If missing_information exists, explain why the claim cannot currently
  be assessed.
- For non-verifiable claims, explain how the next reasoning should name
  one explicit parameter, one measurable metric, and a direction.
- Do not invent evidence.
- Do not rewrite the original claim as if it were established fact.
- Generate guidance that can improve reasoning in the next iteration.

Use evidence_state according to these rules:

- supported:
  The claimed metric direction is observed in an exact
  parameter-isolated runtime comparison.

- contradicted:
  The observed metric direction contradicts the original claim.

- observed_only:
  The claimed trend was observed, but multiple parameters changed
  simultaneously, so the effect cannot be attributed solely to the
  claimed parameter.

- non_verifiable:
  The claim cannot be evaluated because it lacks a measurable metric,
  parameter direction, or expected metric direction.

- insufficient_evidence:
  The claim is measurable, but no suitable runtime comparison currently
  exists.

Return JSON only with this exact structure:

{
  "method": "Reflexion-inspired",
  "iteration_reflection": {
    "summary": "...",
    "reasoning_strengths": [
      "..."
    ],
    "reasoning_issues": [
      "..."
    ],
    "guidance_for_next_iteration": [
      "..."
    ]
  },
  "claim_reflections": [
    {
      "claim_text": "...",
      "evidence_state": "supported/contradicted/observed_only/non_verifiable/insufficient_evidence",
      "runtime_interpretation": "...",
      "reflection": "...",
      "guidance": "...",
      "source_iterations": [
        1
      ]
    }
  ]
}

Input:
<ITERATION_SPECIFIC_REFLECTION_INPUT>
\end{lstlisting}

The prompt requires the LLM to evaluate architectural
claims using only the supplied runtime evidence.
It explicitly prohibits treating previous LLM-generated
statements as ground truth or inferring isolated
parameter effects from confounded observations.

The resulting JSON response contains an
\texttt{iteration\_reflection} object and a
\texttt{claim\_reflections} array. The former
summarizes broader reasoning strengths, issues,
and guidance, while the latter provides
claim-specific evidence assessments and
corrective recommendations.

These outputs are subsequently preserved in
Dynamic Memory and made available during
future architectural candidate generation.

The corresponding recorded Reflection outputs
and their reuse in subsequent reasoning are
illustrated in the V1 case study in
Section~\ref{app:p4_v1_case_study}.

% Insert the verified Reflection prompt here.


\subsection{ExpeL-Based Experiential Learning Prompt (V2)}
\label{app:p4_expel_prompt}

V2 extends the evidence-grounded Reflection mechanism
introduced in V1 by incorporating ExpeL-inspired
experiential learning.

While Reflection identifies weaknesses in previous
architectural reasoning and generates corrective
feedback, ExpeL transforms runtime observations,
claim evidence, and reasoning feedback into reusable
experiential lessons for subsequent DSE iterations.

The Reflection mechanism remains part of V2.
Therefore, the Reflection prompt is not repeated
here; only the additional ExpeL-specific instructions
are presented.

The ExpeL prompt defines four lesson types:
\texttt{observed\_pattern}, \texttt{causal\_warning},
\texttt{reasoning\_policy}, and \texttt{evidence\_gap}.

Listing~\ref{lst:p4_expel_prompt} presents an
abridged version of the ExpeL prompt used in V2.
The excerpt preserves the principal learning
instructions, evidence-authority requirements,
and structured output fields.

\begin{lstlisting}[
    style=remodeljson,
    caption={Abridged ExpeL-specific prompt
    used for experiential learning in V2.},
    label={lst:p4_expel_prompt},
    breaklines=true,
    basicstyle=\ttfamily\footnotesize
]
EXPEL-INSPIRED EXPERIENTIAL LEARNING

Convert each completed DSE iteration into reusable
experiential knowledge derived from the agent's
own runtime simulation history.

Do not merely restate the original extracted claims.

Generate lessons using the following priority:

PRIORITY 1 - OBSERVED RUNTIME PATTERNS

For claims with supporting, contradicting, or
unchanged evidence, generate an observed_pattern.

Include:
- Parameter name and old/new values.
- Metric name and old/new measured values.
- Metric percentage change, when available.
- Other architectural parameters that changed.
- Exact parameter-isolation status.
- Source comparison and current iteration.

If multiple parameters changed, describe the
result as an observed multi-parameter transition,
not as an isolated causal effect.

PRIORITY 2 - CAUSAL WARNINGS

Generate a causal_warning when:
- Multiple architectural parameters changed.
- Evidence is mixed or not parameter-isolated.
- The reasoning attributes a metric change
  to one parameter despite confounding.

Identify the confounding parameters explicitly.

PRIORITY 3 - REASONING POLICIES

Generate a reasoning_policy when:
- Claims are non-verifiable.
- Metrics or expected directions are unclear.
- Hypotheses are presented as observed facts.
- Claims refer to parameters that did not change.

Provide concrete guidance for future reasoning.

PRIORITY 4 - EVIDENCE GAPS

Generate an evidence_gap when:
- No suitable runtime comparison exists.
- The relevant parameter was not changed.
- The metric was unavailable.
- No exact parameter-isolated comparison exists.

Identify the missing evidence and recommend
a useful future controlled comparison.

EVIDENCE AUTHORITY

1. Measured simulation results are authoritative.
2. Exact parameter-isolated evidence is strongest.
3. Multi-parameter observations are not causal proof.
4. Previous claims, reflections, and lessons
   are not empirical facts.
5. Do not use the offline database, optimal design,
   closeness, or external architectural knowledge.
6. Preserve uncertainty.
7. Do not generalize a single observation as
   universally applicable.

Return JSON only:

{
  "method": "ExpeL-inspired",
  "experience_summary": "...",
  "lessons": [
    {
      "lesson_type":
        "observed_pattern/causal_warning/
         reasoning_policy/evidence_gap",
      "lesson": "...",
      "applies_to_parameters": ["..."],
      "applies_to_metrics": ["..."],
      "source_claims": ["..."],
      "source_iterations": [1],
      "evidence_category": "...",
      "exact_parameter_isolation": true,
      "confounding_parameters": ["..."],
      "recommended_future_comparison": "..."
    }
  ]
}

ADDITIONAL REQUIREMENTS

- List all confounding parameters for causal warnings.
- Specify missing controlled comparisons for evidence gaps.
- Avoid duplicate lessons.
- Prefer precise lessons over repetitive lessons.

Input:
<ITERATION_SPECIFIC_EXPERIENCE_INPUT>
\end{lstlisting}

The prompt explicitly prevents ExpeL from
converting unsupported architectural hypotheses
into established knowledge. It also requires
each lesson to retain its source evidence and
identify whether the corresponding observations
provide exact parameter isolation.

The resulting structured lessons are preserved
in Dynamic Memory and made available to
subsequent ACG iterations.

The generation and subsequent reuse of these
lessons are illustrated using recorded
experimental outputs in
Section~\ref{app:p4_v2_case_study}.



\subsection{Runtime Evidence Graph-Guided Prompt (V3)}
\label{app:p4_graph_prompt}

V3 extends the Reflection and ExpeL mechanisms
by introducing a Runtime Evidence Graph that
organizes accumulated architectural claims,
simulation observations, and experiential
feedback into explicit parameter--metric
relationships.

Unlike Reflection and ExpeL, graph construction
and relationship ranking are performed through
deterministic processing rather than an
additional LLM invocation.

The LLM-based component of V3 is the
graph-informed ACG stage, which receives
a compact representation of accumulated
relationships and their evidence assessments.

The Reflection and ExpeL prompts remain
as described in
Sections~\ref{app:p4_reflection_prompt}
and~\ref{app:p4_expel_prompt}.
Therefore, this subsection presents only
the additional graph-specific instructions
and refinement-context fields introduced
in V3.

\paragraph{Compact Graph Context Supplied to ACG.}

The complete Runtime Evidence Graph remains
available internally. However, ACG receives
a compact relationship-level representation
rather than the full graph structure.

The graph-specific refinement context
contains three principal components:

\begin{enumerate}

\item \texttt{runtime\_graph\_relationship\_view}:
Provides the complete set of accumulated
parameter--metric relationships in compact
form, including their evidence states,
causal-strength assessments, and
source iterations.

\item \texttt{runtime\_relationship\_ranking}:
Provides deterministic evidence-aware
ranking and recommended actions for
the recorded relationships.

\item \texttt{relationship\_evidence\_details}:
Provides additional evidence-quality
information, iteration outcomes,
and reasoning policies associated
with individual relationships.

\end{enumerate}

The compact representation excludes
raw claim prose, Reflection text,
ExpeL lesson text, and graph edges.
These remain available through the
internally maintained graph.

\paragraph{Graph-Specific ACG Instructions.}

Listing~\ref{lst:p4_graph_prompt}
presents the graph-specific instructions
extracted from the V3 ACG input
construction in \texttt{app/agent.py}.

The listing omits general candidate-generation
instructions already discussed in the
REMODEL-V framework and preserves the
graph-specific guidance used in V3.

\begin{lstlisting}[
    style=remodeljson,
    caption={Graph-specific ACG instructions
    used in V3.},
    label={lst:p4_graph_prompt},
    breaklines=true,
    basicstyle=\ttfamily\footnotesize
]
Use runtime_relationship_ranking as the primary ordered
evidence summary when planning new candidates.

Use runtime_graph_relationship_view to inspect every
Runtime Evidence Graph relationship, including lower-ranked
or currently unranked relationships.

Use relationship_evidence_details when evidence quality,
iteration outcomes, isolation quality, or uncertainty
must be examined.

Use relationship_state, causal_strength, evidence_counts,
iteration_outcomes, and source_iterations when evaluating
prior claims.

Treat runtime graph relationship states and measured
iteration outcomes as higher-authority information
than generated guidance.

When prior runtime evidence was not exactly
parameter-isolated, describe it only as an observed
multi-parameter relationship, not as a causal effect.

Do not repeat a claim that prior runtime evidence
contradicted unless you clearly identify it as
a new hypothesis requiring another execution.

When prior Runtime Evidence Graph evidence influences
a design decision, identify the source iteration
in the retrospect or analysis.

Clearly distinguish prior measured runtime
observations from new hypotheses.

Do not use raw history or invent previous iterations.

MANDATORY GRAPH-GROUNDING RULE:

For every changed parameter and discussed metric,
inspect runtime_graph_relationship_view.relationships
for a relationship whose parameter,
parameter_direction, and metric match the
proposed change.

In the analysis, report its node_id,
relationship_state, causal_strength,
and source_iterations.

If no matching relationship exists, label
the expected effect as a new unverified hypothesis.

If evidence is mixed, contradicted, unchanged,
or insufficient, state the uncertainty
and do not claim improvement.

Treat cost feasibility for an unexecuted
candidate only as an expectation requiring execution.
\end{lstlisting}

The mandatory graph-grounding rule requires
the LLM to connect its architectural predictions
to recorded graph relationships whenever
matching evidence is available.

When no matching relationship exists,
the proposed effect must be identified
as a new unverified hypothesis.

Furthermore, the prompt requires the
LLM to distinguish observed multi-parameter
transitions from causal effects and to
acknowledge contradictory, mixed, or
insufficient evidence.

Thus, V3 introduces explicit evidence-use
requirements during candidate generation,
rather than relying solely on the LLM
to independently interpret accumulated
textual experiences.

The corresponding graph structure and
recorded execution behavior are illustrated
in Section~\ref{app:p4_v3_case_study}.


\section{Recorded Execution Case Studies}
\label{app:p4_case_studies}

\subsection{Evidence-Grounded Reflection (V1)}
\label{app:p4_v1_case_study}
\subsubsection{Step 1: Architectural Reasoning Generated by ACG}

As described in Appendix~A, Architectural Candidate
Generation (ACG) uses an LLM to propose architectural
configurations and provide structured reasoning
for the expected effects of parameter changes.

During Iteration 4, ACG proposed a configuration
containing 192 reorder-buffer (ROB) entries.

The recorded ACG reasoning included the following
prediction:

\begin{quote}
\emph{``Setting rob\_entries to 192 may increase ipc.''}
\end{quote}

This prediction expresses the LLM's expectation
that increasing ROB capacity may improve
instructions per cycle (IPC).

At this stage, the statement is an architectural
hypothesis generated by the LLM. Its correctness
has not yet been established by runtime evidence.

The selected configuration was subsequently
evaluated using the architectural execution
backend, producing an IPC of approximately
1.148105.


% =====================================================
% STEP 2: CLAIM EXTRACTION
% =====================================================

\subsubsection{Step 2: Extracting the Architectural Claim}

After execution, the Claim Extractor processes
the reasoning associated with the selected
architectural candidate.

The Claim Extractor uses an LLM with a
constrained extraction prompt to identify
atomic architectural predictions from
the original ACG reasoning.

Its responsibility is to identify what the
LLM predicted, not to determine whether
the prediction is correct.

For the recorded ROB--IPC reasoning, the
Claim Extractor produces the following
structured representation.

\begin{lstlisting}[
    style=remodeljson,
    caption={Recorded architectural claim extracted
    from the Iteration 4 ACG reasoning.},
    label={lst:app_v1_extracted_claim}
]
{
  "claim_text":
    "Increasing rob_entries may increase ipc.",
  "claim_type": "verifiable_claim",
  "reasoning_type": "parameter_change",
  "design_parameter_changed": "rob_entries",
  "parameter_direction": "increase",
  "new_value": "192",
  "metric": "ipc",
  "expected_direction": "increase",
  "source_section": "analysis.steps[5]"
}
\end{lstlisting}

The extracted representation identifies four
important elements:

\begin{enumerate}
    \item \textbf{Architectural parameter:}
    \texttt{rob\_entries}, representing ROB capacity.

    \item \textbf{Parameter-change direction:}
    \texttt{increase}, indicating a proposed
    increase in ROB entries.

    \item \textbf{Target metric:}
    \texttt{ipc}, representing instructions
    per cycle.

    \item \textbf{Expected metric direction:}
    \texttt{increase}, indicating that the
    LLM expects IPC to improve.
\end{enumerate}

Because the claim identifies a valid
architectural parameter, a measurable metric,
and explicit expected directions, it is
classified as verifiable.

Importantly, this classification does not
mean that the claim is correct. It means
that the prediction contains sufficient
information to be evaluated against
architectural measurements.


% =====================================================
% STEP 3: RUNTIME EVIDENCE
% =====================================================

\subsubsection{Step 3: Collecting Runtime Evidence}

Once the architectural claim has been
extracted, the Claim-Evidence Builder
examines the recorded simulation outcomes
to identify observations relevant to
the predicted parameter--metric relationship.

In this example, the prediction is:

\begin{equation}
\mathrm{ROB\ Entries}\uparrow
\;\rightarrow\;
\mathrm{IPC}\uparrow.
\end{equation}

The design executed during Iteration 4
contains 192 ROB entries and achieves
an IPC of 1.148105.

The Claim-Evidence Builder compares this
design against three previously executed
configurations, each containing 128 ROB
entries.

Table~\ref{tab:app_v1_rob_evidence}
summarizes the recorded measurements.

\begin{table}[htbp]
\centering
\caption{Runtime evidence collected for the
ROB--IPC architectural prediction.}
\label{tab:app_v1_rob_evidence}
\renewcommand{\arraystretch}{1.15}
\small
\setlength{\tabcolsep}{4pt}
\resizebox{\columnwidth}{!}{%
\begin{tabular}{lccccc}
\toprule
\textbf{Comparison} &
\textbf{Previous ROB} &
\textbf{Current ROB} &
\textbf{Previous IPC} &
\textbf{Current IPC} &
\textbf{Evidence} \\
\midrule
Iteration 1 vs. 4 &
128 & 192 & 1.144713 & 1.148105 &
Supporting \\

Iteration 2 vs. 4 &
128 & 192 & 1.151127 & 1.148105 &
Contradicting \\

Iteration 3 vs. 4 &
128 & 192 & 1.151127 & 1.148105 &
Contradicting \\
\bottomrule
\end{tabular}%
}
\end{table}

The recorded comparisons are interpreted
as follows:

\begin{enumerate}

\item \textbf{Iteration 1 compared with Iteration 4:}

ROB capacity increases from 128 to 192,
while IPC increases from approximately
1.144713 to 1.148105.

Because both changes follow the predicted
directions, this comparison is classified
as supporting evidence.

\item \textbf{Iteration 2 compared with Iteration 4:}

ROB capacity again increases from 128 to
192, but IPC decreases from approximately
1.151127 to 1.148105.

The observed IPC direction is opposite
to the prediction. Therefore, this
comparison is classified as contradicting
evidence.

\item \textbf{Iteration 3 compared with Iteration 4:}

The same ROB increase is observed, while
IPC decreases from approximately 1.151127
to 1.148105.

This comparison is also classified as
contradicting evidence.

\end{enumerate}

Overall, the Claim-Evidence Builder records
one supporting comparison and two
contradicting comparisons.

However, all three comparisons involve
changes to additional architectural
parameters. None provides exact
parameter isolation.

Consequently, the observed IPC differences
cannot be attributed exclusively to ROB
capacity. The evidence demonstrates that
the predicted relationship is not
consistently supported by the available
runtime comparisons.

The collected evidence is then supplied
to the Reflection mechanism, which
examines the discrepancy between the
LLM-generated prediction and the
observed architectural behavior.


% =====================================================
% STEP 4: EVIDENCE-GROUNDED REFLECTION
% =====================================================

\subsubsection{Step 4: Generating Evidence-Grounded Reflection}

After collecting runtime evidence, V1 invokes
the Reflexion-inspired mechanism to examine
the architectural predictions and identify
reasoning weaknesses.

Unlike the Claim Extractor, which identifies
testable predictions, and the Claim-Evidence
Builder, which collects relevant observations,
Reflection interprets the available evidence
and generates corrective guidance for
subsequent architectural reasoning.

In the selected example, the Claim-Evidence
Builder identified one supporting and two
contradicting comparisons for the predicted
ROB--IPC relationship.

The Reflection mechanism receives this
evidence and generates both claim-specific
feedback and an iteration-level summary.

\paragraph{Recorded Claim-Level Reflection}

Listing~\ref{lst:app_v1_reflection}
presents the recorded Reflection generated
for the ROB--IPC prediction during
Iteration 4.

\begin{lstlisting}[
    style=remodeljson,
    caption={Recorded claim-level Reflection
    for the ROB--IPC prediction.},
    label={lst:app_v1_reflection}
]
{
  "claim_text":
    "Increasing rob_entries may increase ipc.",

  "evidence_state": "contradicted",

  "runtime_interpretation":
    "IPC decreased from 1.15112703986 to
     1.14810454489 despite increasing
     rob_entries from 128 to 192.",

  "reflection":
    "The claim is contradicted by the
     observed decrease in IPC.",

  "guidance":
    "Reassess the relationship between
     rob_entries and IPC, considering
     other influencing parameters."
}
\end{lstlisting}

The recorded Reflection contains four
important components:

\begin{enumerate}

\item \textbf{Architectural claim:}
The \texttt{claim\_text} field identifies
the prediction being evaluated.

\item \textbf{Evidence assessment:}
The \texttt{evidence\_state} field records
the Reflection mechanism's classification
of the claim as \texttt{contradicted}.

\item \textbf{Runtime interpretation:}
The \texttt{runtime\_interpretation} field
describes the observed decrease in IPC
despite the increase in ROB capacity.

\item \textbf{Corrective guidance:}
The \texttt{guidance} field recommends
reassessing the predicted ROB--IPC
relationship while considering other
architectural parameters.

\end{enumerate}

Importantly, the recorded claim-level
Reflection emphasizes the contradictory
comparison, although the complete evidence
record also contains one supporting
observation.

Therefore, the Reflection classification
should not be interpreted as demonstrating
that every available runtime comparison
contradicts the prediction.

Furthermore, because the comparisons
involve changes to multiple architectural
parameters, the evidence does not establish
a causal relationship between ROB capacity
and IPC.


\paragraph{Iteration-Level Reflection}

In addition to evaluating individual claims,
the Reflection mechanism generates an
iteration-level assessment summarizing
broader reasoning strengths and weaknesses.


The following excerpt is taken directly from
the recorded Reflection response. Only the
reasoning issues and guidance fields are shown;
the summary and reasoning-strength fields are
omitted for brevity.


Listing~\ref{lst:app_v1_iteration_reflection}
presents selected fields from the recorded
Iteration 4 Reflection.

\begin{lstlisting}[
    style=remodeljson,
    caption={Selected fields from the recorded
    iteration-level Reflection generated after
    Iteration 4 in V1.},   
    label={lst:app_v1_iteration_reflection}
]
{
  "reasoning_issues": [
    "Some claims lack measurable metrics,
     making them non-verifiable.",

    "Contradictions in claims regarding
     power consumption and IPC indicate
     potential misinterpretations of
     parameter effects."
  ],

  "guidance_for_next_iteration": [
    "Ensure all claims specify measurable
     metrics and expected directions.",

    "Focus on isolating parameters in
     future designs to clarify causality.",

    "Review contradictory claims to
     refine understanding of parameter
     interactions."
  ]
}
\end{lstlisting}

The iteration-level Reflection identifies
three areas requiring improvement.

First, architectural predictions should
specify measurable metrics and explicit
expected directions to enable verification.

Second, parameter isolation should be
considered when interpreting architectural
effects, since multiple simultaneous changes
can introduce confounding.

Third, previously contradicted predictions
should be reconsidered before being reused
in subsequent candidate generation.

These recommendations demonstrate how V1
converts runtime observations into
structured corrective feedback.

However, generating such feedback does
not necessarily guarantee that the LLM
will apply it correctly in later iterations.


% Insert the completed V1 case study here.



% =====================================================
% STEP 5: DYNAMIC MEMORY AND FEEDBACK REUSE
% =====================================================

\subsubsection{Step 5: Preserving and Reusing Reflection Feedback}

After Reflection generation, V1 stores the
executed architectural configuration,
measured simulation outcomes, extracted
claims, associated runtime evidence, and
Reflection feedback in Dynamic Memory.

At the beginning of Iteration 5, the
recorded refinement context contains four
previous experiences and four Reflection
records.

These records are made available to ACG
so that subsequent architectural proposals
can incorporate lessons from earlier
execution outcomes.

\paragraph{Recorded Subsequent ACG Reasoning}

Despite the availability of the preceding
Reflection, the Iteration 5 ACG reasoning
again associates 192 ROB entries with
an expected IPC improvement.

Listing~\ref{lst:app_v1_next_acg}
presents the recorded reasoning excerpt.

\begin{lstlisting}[
    style=remodeljson,
    caption={Recorded Iteration 5 ACG reasoning
    following evidence-grounded Reflection.},
    label={lst:app_v1_next_acg}
]
{
  "thinking":
    "Maintaining the ROB entries at 192
     can support instruction throughput
     effectively without violating the
     cost constraints, thus preserving
     performance.",

  "step_conclusion":
    "Setting rob_entries to 192 may
     increase ipc."
}
\end{lstlisting}

The generated reasoning again predicts
a favorable relationship between ROB
capacity and IPC.

However, it does not explicitly acknowledge
the contradictory observation identified
during Iteration 4 or explain why the
previously recorded uncertainty should
no longer apply.

This observation illustrates that making
corrective feedback available in Dynamic
Memory does not guarantee its explicit
application in subsequent architectural
reasoning.

It does not establish that all Reflection
feedback was ignored. Rather, it identifies
a specific instance in which the LLM
repeated a favorable architectural
prediction without addressing the earlier
contradictory evidence.


\paragraph{Interpretation of the V1 Case Study}

The recorded execution demonstrates the
operation and limitations of evidence-grounded
Reflection.

First, V1 extracts architectural predictions
and connects them to observed simulation
evidence.

Second, Reflection transforms the identified
reasoning weaknesses into claim-specific
and iteration-level corrective guidance.

Third, Dynamic Memory preserves this feedback
and makes it available during subsequent
candidate generation.

However, the Iteration 5 reasoning demonstrates
that the availability of corrective feedback
does not ensure its consistent application.

This case study therefore illustrates
the distinction between generating corrective
knowledge and effectively reusing that
knowledge during iterative DSE.

The observed limitation motivates the
investigation of additional mechanisms,
including ExpeL-based experiential learning
in V2 and structured evidence organization
through the Runtime Evidence Graph in V3.




\subsection{ExpeL-Based Experiential Learning (V2)}
\label{app:p4_v2_case_study}

This case study illustrates how V2 extends
evidence-grounded Reflection by introducing
ExpeL-inspired experiential learning.

While V1 generates corrective feedback from
runtime evidence, V2 additionally converts
execution observations and reasoning weaknesses
into reusable lessons stored in Dynamic Memory.


Two recorded executions from the Matrix
Multiplication scenario under the Wide constraint
regime are examined. The primary walkthrough
uses \texttt{run\_20260818\_234443} to illustrate
the generation, preservation, and availability
of ExpeL lessons, followed by an examination
of subsequent ACG reasoning during Iterations
3 and 4.

A complementary execution,
\texttt{run\_20260818\_224903}, illustrates
evidence-consistent reasoning in which a
subsequent ACG decision aligns with previously
recorded ExpeL observations.

Together, these examples demonstrate both
the capabilities and limitations of
ExpeL-inspired experiential learning.



% =====================================================
% STEP 1: ARCHITECTURAL PREDICTION
% =====================================================

\subsubsection{Step 1: Architectural Prediction and Runtime Observation}

During Iteration 3, ACG generated the following
architectural prediction:

\begin{quote}
\emph{``Increasing l1d\_size may increase hit\_rate.''}
\end{quote}

The Claim Extractor identified this statement
as a verifiable prediction involving an increase
in L1D cache size and an expected increase in
cache hit rate.

Following execution, the Claim-Evidence Builder
examined the corresponding runtime observations.

Table~\ref{tab:app_v2_l1d_evidence}
summarizes the recorded comparison between
Iterations 2 and 3.

\begin{table}[htbp]
\centering
\caption{Recorded L1D cache size and hit-rate
measurements in the V2 case study.}
\label{tab:app_v2_l1d_evidence}
\renewcommand{\arraystretch}{1.15}
\small
\begin{tabular}{lcc}
\toprule
\textbf{Measurement} &
\textbf{Iteration 2} &
\textbf{Iteration 3} \\
\midrule
L1D cache size & 8 kB & 32 kB \\
Hit rate & 0.999986879309 & 0.999986879137 \\
\bottomrule
\end{tabular}
\end{table}

The observed hit rate decreased slightly,
despite the increase in L1D cache size.
Thus, the observed metric direction does
not agree with the original prediction.

However, the comparison also involves
changes in L1D MSHRs, L2 associativity,
L2 cache size, load queue entries, and
ROB entries.

Therefore, the observed hit-rate difference
cannot be attributed exclusively to the
change in L1D cache size.

This distinction is important because
ExpeL must preserve the observed runtime
outcome without converting it into an
unsupported causal conclusion.


% =====================================================
% STEP 2: REFLECTION
% =====================================================

\subsubsection{Step 2: Reflection Identifies the Reasoning Weakness}

As in V1, the Reflection mechanism examines
the extracted claims and associated runtime
evidence before generating corrective feedback.

For Iteration 3, the recorded Reflection
identified contradictions involving
\texttt{l1d\_size} and \texttt{l1d\_mshrs}.

The iteration-level guidance included
the following recommendations:

\begin{quote}
\emph{``Focus on isolating parameters more
effectively to establish clearer causal
relationships.''}
\end{quote}

\begin{quote}
\emph{``Ensure that all claims specify
measurable metrics and expected directions
to enhance verifiability.''}
\end{quote}

These recommendations identify reasoning
weaknesses and suggest how subsequent
architectural reasoning should improve.

V2 then introduces an additional ExpeL stage
to convert these observations and corrective
insights into reusable experiential lessons.


% =====================================================
% STEP 3: EXPEL LESSON GENERATION
% =====================================================

\subsubsection{Step 3: ExpeL Generates Reusable Lessons}

The ExpeL stage receives runtime observations,
claim evidence, and Reflection feedback and
generates structured experiential lessons.

As described in
Section~\ref{app:p4_expel_prompt},
the ExpeL prompt supports four lesson types:
\texttt{observed\_pattern},
\texttt{causal\_warning},
\texttt{reasoning\_policy}, and
\texttt{evidence\_gap}.

In the selected example, ExpeL generates
an observed-pattern lesson describing
the L1D-size and hit-rate comparison.

Listing~\ref{lst:app_v2_observed_pattern}
presents selected fields from the recorded
lesson.

\begin{lstlisting}[
    style=remodeljson,
    caption={Selected fields from the recorded
    ExpeL observed-pattern lesson in Iteration 3.},
    label={lst:app_v2_observed_pattern},
    breaklines=true
]
{
  "lesson_type": "observed_pattern",
  "source_iterations": [2, 3],
  "evidence_category": "contradicting",
  "exact_parameter_isolation": false,
  "applies_to_parameters": [
    "l1d_size",
    "l1d_mshrs",
    "l2_assoc",
    "l2_size",
    "lq_entries",
    "rob_entries"
  ],
  "applies_to_metrics": ["hit_rate"],
  "confounding_parameters": [
    "l1d_mshrs",
    "l2_assoc",
    "l2_size",
    "lq_entries",
    "rob_entries"
  ]
}
\end{lstlisting}

The recorded lesson also describes the
measured decrease in hit rate from
0.999986879309 to 0.999986879137
as L1D cache size increased from
8 kB to 32 kB.

Unlike a general textual recommendation,
the ExpeL lesson retains the source
iterations, affected architectural
parameters, target metric, and evidence
classification.

The field
\texttt{exact\_parameter\_isolation=false}
indicates that the comparison involved
simultaneous changes to other parameters.

ExpeL also generates a causal-warning
lesson for the same execution.

\begin{lstlisting}[
    style=remodeljson,
    caption={Recorded ExpeL causal warning
    associated with the Iteration 3 observations.},
    label={lst:app_v2_causal_warning},
    breaklines=true
]
{
  "lesson_type": "causal_warning",
  "lesson": "The observed metric changes for l1d_size and l1d_mshrs cannot be causally assigned due to simultaneous changes in multiple parameters, including l2_assoc, l2_size, lq_entries, and rob_entries.",
  "source_iterations": [2, 3],
  "exact_parameter_isolation": false
}
\end{lstlisting}

This warning explicitly explicitly cautions against the
observed metric changes from being
interpreted as isolated causal effects
of L1D cache size or MSHR capacity.

Together, the observed-pattern and
causal-warning lessons demonstrate
the additional functionality introduced
by ExpeL: converting execution outcomes
into structured, reusable experiential
knowledge while retaining evidence
limitations.


% =====================================================
% STEP 4: DYNAMIC MEMORY
% =====================================================

\subsubsection{Step 4: Preserving ExpeL Lessons in Dynamic Memory}

After lesson generation, V2 stores the
ExpeL outputs in Dynamic Memory alongside
the executed configurations, measured
simulation outcomes, extracted claims,
and Reflection feedback.

At the beginning of Iteration 4, the
recorded refinement context contained
three previous experiences, three
Reflection records, and fifteen
accumulated ExpeL lessons.

Table~\ref{tab:app_v2_memory_growth}
summarizes the recorded memory availability
across subsequent ACG iterations.

\begin{table}[htbp]
\centering
\caption{Accumulated Reflection and ExpeL
records available to subsequent ACG iterations.}
\label{tab:app_v2_memory_growth}
\renewcommand{\arraystretch}{1.15}
\small
\begin{tabular}{lccc}
\toprule
\textbf{ACG Iteration} &
\textbf{Experiences} &
\textbf{Reflections} &
\textbf{ExpeL Lessons} \\
\midrule
2 & 1 & 1 & 4 \\
3 & 2 & 2 & 9 \\
4 & 3 & 3 & 15 \\
5 & 4 & 4 & 21 \\
\bottomrule
\end{tabular}
\end{table}

These records confirm that ExpeL lessons
accumulate across completed DSE iterations
and remain available to subsequent
candidate generation.

However, the presence of lessons in
Dynamic Memory does not necessarily
demonstrate that the LLM incorporates
them correctly into its reasoning.



% =====================================================
% STEP 5: SUBSEQUENT ACG
% =====================================================

\subsubsection{Step 5: Examining Subsequent ACG Reasoning}

To examine how the stored lessons are
used, the reasoning generated during
Iteration 4 is inspected.

The ACG refinement guidance explicitly
instructed the LLM to review previous
claim evidence, Reflection feedback,
and experiential lessons.

Nevertheless, the selected ACG reasoning
again proposed increasing L1D cache size
as a means of improving hit rate.

The recorded reasoning stated:

\begin{quote}
\emph{``Increasing l1d\_size may increase
hit\_rate.''}
\end{quote}

The associated reasoning proposed
increasing L1D cache size to 64 kB,
suggesting that this configuration
could enhance hit rate without
significantly affecting cost metrics.

However, the generated reasoning did
not explicitly acknowledge the earlier
contradictory observation or the
causal-warning lesson preserved
in Dynamic Memory.

This example demonstrates that although
ExpeL successfully generates and
forwards reusable lessons, their
availability does not guarantee
consistent application during subsequent
architectural reasoning.

Importantly, this observation does not
establish that the LLM ignored all
previous experiences or lessons.
It identifies a specific instance
in which the generated reasoning
repeated a favorable prediction
without addressing previously
recorded uncertainty.


\subsubsection{Example of Evidence-Consistent ExpeL Lesson Reuse}
\label{app:p4_v2_successful_reuse}

To complement the preceding example, an additional
V2 execution is examined to illustrate a case in
which subsequent ACG reasoning is consistent with
previously generated ExpeL lessons.

The example is taken from the Matrix Multiplication
scenario under the Wide constraint regime,
identified by \texttt{run\_20260818\_224903}.

\paragraph{Recorded ExpeL Lessons.}

During the earlier DSE iterations, ExpeL generated
observed-pattern lessons describing architectural
transitions in which reductions in L1D cache size
were accompanied by lower measured dynamic power.

Table~\ref{tab:app_v2_positive_expel}
summarizes two recorded observations.

\begin{table}[htbp]
\centering
\caption{Recorded ExpeL observations involving
L1D cache size and dynamic power.}
\label{tab:app_v2_positive_expel}
\renewcommand{\arraystretch}{1.15}
\small
\begin{tabular}{lcc}
\toprule
\textbf{Measurement} &
\textbf{Observation 1} &
\textbf{Observation 2} \\
\midrule
L1D size &
32 $\rightarrow$ 16 kB &
32 $\rightarrow$ 8 kB \\

Dynamic power &
2.54076 $\rightarrow$ 2.23183 &
2.54076 $\rightarrow$ 2.12807 \\

Exact parameter isolation &
No & No \\
\bottomrule
\end{tabular}
\end{table}

The first observation was recorded in an
ExpeL lesson generated after Iteration 2,
while the second was recorded after Iteration 3.

Both observations describe decreases in
dynamic power during architectural transitions
that included reductions in L1D cache size.

However, other parameters also changed.
The first observation involved changes in
L1D MSHRs, while the second involved changes
in L2 cache size.

Consequently, these observations do not
establish that reducing L1D cache size alone
caused the measured power reductions.

ExpeL preserved this limitation through
\texttt{exact\_parameter\_isolation=false}
and generated corresponding causal-warning
lessons.

\paragraph{Availability in Dynamic Memory.}

At the beginning of Iteration 4, the recorded
Dynamic Memory input contained three previous
experiences and twenty accumulated ExpeL
lessons.

These records included the earlier observed
power reductions and their associated
causal warnings.

Thus, the ExpeL lessons were available to
ACG before the next candidate was generated.

\paragraph{Subsequent ACG Reasoning.}

During Iteration 4, ACG proposed reducing
L1D cache size to 16 kB as part of a
cost-oriented architectural configuration.

The recorded reasoning included:

\begin{quote}
\emph{``Reducing l1d\_size to '16kB' may decrease
total\_power due to a smaller cache size,
potentially lowering dynamic power consumption.''}
\end{quote}

The final explanation further stated:

\begin{quote}
\emph{``This candidate configuration attempts
to lower the L1D size to reduce power metrics,
which has previously resulted in favorable
outcomes for cost metrics based on observed
evidence.''}
\end{quote}

These statements are consistent with the
earlier ExpeL lessons documenting observed
power reductions during transitions involving
smaller L1D cache configurations.

Importantly, the LLM explicitly references
previously observed evidence when explaining
its design choice.

Nevertheless, the generated reasoning does
not explicitly acknowledge the confounding
parameters identified in the earlier
causal-warning lessons.

Therefore, this example provides evidence of
subsequent reasoning that is consistent with
previously recorded ExpeL observations. However, it does not demonstrate complete application
of the associated causal warnings.

\paragraph{Interpretation.}

This example illustrates that ExpeL can
produce reusable observations that are
consistent with subsequent architectural
decisions.

However, the recorded trace does not
establish that ExpeL alone caused the
decision, since ACG also receives other
execution history and feedback through
Dynamic Memory.

Taken together with the preceding
limitation example, these observations
show that experiential lessons can
support subsequent reasoning, but their
application remains dependent on how
the LLM interprets and incorporates
the available textual evidence.


% =====================================================
% CASE STUDY INTERPRETATION
% =====================================================

\subsubsection{Case Study Interpretation}

The two recorded V2 executions illustrate
both the capabilities and limitations of
ExpeL-inspired experiential learning.

First, ExpeL transforms runtime observations
and Reflection feedback into reusable lessons
that retain architectural parameters,
measured outcomes, evidence classifications,
and source iterations.

Second, the mechanism preserves important
evidence limitations by identifying
multi-parameter observations and generating
causal warnings when parameter effects
cannot be isolated.

Third, the complementary execution
demonstrates that subsequent ACG reasoning
can align with previously recorded
experiential observations and explicitly
refer to prior evidence when motivating
architectural decisions.

However, the primary execution shows that
the availability of corrective lessons
does not guarantee their explicit application.
In particular, ACG may repeat a favorable
architectural prediction without acknowledging
previously recorded contradictory observations
or causal warnings.

Together, these examples demonstrate that
ExpeL provides a mechanism for generating,
preserving, and forwarding experiential
knowledge, but the effective use of this
knowledge remains dependent on the LLM's
interpretation of accumulated textual feedback.

This limitation motivates the structured
evidence representation introduced in V3,
where architectural claims and runtime
observations are organized into explicit
parameter--metric relationships with
evidence-aware states and deterministic
ranking.


\subsection{Runtime Evidence Graph (V3)}
\label{app:p4_v3_case_study}

This case study illustrates how V3 extends the
Reflection and ExpeL mechanisms introduced in V1
and V2 by organizing accumulated architectural
evidence into a structured Runtime Evidence Graph.

While V2 stores executed designs, simulation
results, extracted claims, Reflection feedback,
and experiential lessons in iteration-based
Dynamic Memory, V3 represents these information
types as interconnected graph nodes and edges.

The purpose of this representation is to make
architectural relationships and their supporting
evidence explicitly identifiable, assessable,
and reusable during subsequent DSE iterations.

The example is taken from a recorded Matrix
Multiplication experiment under the Wide
constraint regime, identified by
\texttt{run\_20260731\_195619}.

The graph snapshot after Iteration 3 is used
to illustrate how the framework connects
architectural claims, runtime evidence, and
experiential lessons into shared
parameter--metric relationships.


% =====================================================
% STEP 1: GRAPH ORGANIZATION
% =====================================================

\subsubsection{Step 1: Organizing Runtime Evidence into a Graph}

After each DSE iteration, V3 updates the
Runtime Evidence Graph using the executed
architectural configuration, measured
simulation outcomes, extracted claims,
associated runtime evidence, Reflection
feedback, and ExpeL lessons.

Unlike an iteration-based textual memory,
the graph explicitly represents these
information types as nodes connected
through typed edges.

The full graph contains several node types,
including iteration, design, parameter,
metric, claim, evidence, Reflection,
lesson, and parameter--metric relationship
nodes.

For example, an iteration node is linked
to the architectural design executed
during that iteration. The design node
is connected to its parameter values
and measured performance and cost metrics.

Similarly, an extracted claim is connected
to its referenced architectural parameter,
predicted metric, and associated runtime
evidence.

Claims describing the same parameter
change and expected metric direction
are additionally associated with a
shared parameter--metric relationship.

This organization allows multiple
architectural predictions to contribute
evidence to the same relationship without
discarding their individual source records.


% =====================================================
% STEP 2: RECORDED GRAPH SNAPSHOT
% =====================================================

\subsubsection{Step 2: Examining a Recorded Graph Snapshot}

The selected Runtime Evidence Graph snapshot
was recorded after Iteration 3 of the
Matrix Multiplication experiment.

At this point, the graph contained three
executed iterations, three architectural
designs, twenty-six extracted claims,
twenty-six evidence nodes, and four unique
parameter--metric relationship records.

The graph also contained Reflection and
ExpeL records associated with the completed
iterations.

Table~\ref{tab:app_v3_graph_snapshot}
summarizes the recorded graph structure.

\begin{table}[htbp]
\centering
\caption{Recorded Runtime Evidence Graph
structure after Iteration 3.}
\label{tab:app_v3_graph_snapshot}
\renewcommand{\arraystretch}{1.15}
\small
\begin{tabular}{lc}
\toprule
\textbf{Graph Component} &
\textbf{Recorded Count} \\
\midrule
Iteration nodes & 3 \\
Executed design nodes & 3 \\
Parameter nodes & 6 \\
Metric nodes & 14 \\
Claim nodes & 26 \\
Evidence nodes & 26 \\
Iteration-level Reflection nodes & 3 \\
Claim-level Reflection nodes & 23 \\
ExpeL lesson nodes & 10 \\
Parameter--metric relationship nodes & 4 \\
Graph edges & 330 \\
\bottomrule
\end{tabular}
\end{table}


These measurements illustrate that the
Runtime Evidence Graph retains several
types of information generated during DSE,
rather than storing only architectural
relationships.

Although the complete Runtime Evidence Graph
is maintained internally, forwarding every
graph node and edge to ACG would introduce
unnecessary contextual detail.

Therefore, V3 constructs a compact
relationship-level representation for
subsequent candidate generation.

This representation includes the recorded
parameter--metric relationships, their
evidence states, source iterations,
evidence-quality information, and
evidence-aware rankings.

Consequently, ACG can inspect accumulated
architectural relationships without
receiving the complete graph structure
or the full textual history of previous
experiences.



% =====================================================
% STEP 3: REPRESENTATIVE RELATIONSHIP
% =====================================================

\subsubsection{Step 3: Understanding a Parameter--Metric Relationship}

To illustrate how architectural evidence
is organized, consider the following
relationship recorded in the graph:

\begin{equation}
\mathrm{L1D\ Size}\downarrow
\;\rightarrow\;
\mathrm{Area}\downarrow.
\end{equation}

This relationship represents the prediction
that decreasing L1D cache size may be
associated with a reduction in area.

The corresponding relationship record is
identified as:

\begin{quote}
\texttt{relationship::l1d\_size::decrease::area::decrease}
\end{quote}

In the selected graph snapshot, three
extracted claims from Iterations 2 and 3
are associated with this relationship.

Each claim retains its corresponding
runtime evidence record, while three
ExpeL lesson nodes are also linked
to the relationship.

The relationship therefore acts as a
shared evidence record connecting
architectural predictions that describe
the same expected parameter--metric behavior.


% =====================================================
% STEP 4: GRAPH NODES AND CONNECTIONS
% =====================================================

\subsubsection{Step 4: Understanding Graph Nodes and Connections}

The Runtime Evidence Graph represents
architectural information as interconnected
nodes and typed edges.

Nodes correspond to individual information
records, such as extracted claims, runtime
evidence, Reflection feedback, ExpeL lessons,
and parameter--metric relationships.

Edges define how these records are connected
and how information from different DSE
iterations contributes to accumulated
architectural evidence.

For the recorded L1D-size--area relationship
introduced in Step 3, three extracted claims
from Iterations 2 and 3 are associated with
a common relationship node.

Each claim is connected to its corresponding
runtime evidence record, while three ExpeL
lesson nodes are also linked to the same
relationship.

The recorded graph uses four important
types of connections:

\begin{enumerate}

\item \textbf{Claim to Evidence:}
The \texttt{CLAIM\_EVALUATED\_BY} edge connects
an extracted architectural prediction to
its corresponding runtime evidence record.

\item \textbf{Claim to Relationship:}
The \texttt{CLAIM\_CONTRIBUTES\_TO\_RELATIONSHIP}
edge associates an architectural prediction
with the shared parameter--metric relationship
that it describes.

\item \textbf{Evidence to Relationship:}
The \texttt{EVIDENCE\_UPDATES\_RELATIONSHIP}
edge connects runtime evidence records
to the relationship whose accumulated
evidence assessment they contribute to.

\item \textbf{Lesson to Relationship:}
The \texttt{LESSON\_INFORMS\_RELATIONSHIP}
edge associates ExpeL-generated lessons
with the corresponding architectural
relationship.

\end{enumerate}

In the selected graph snapshot, the
L1D-size--area relationship is connected
to three claim nodes, three evidence
nodes, and three ExpeL lesson nodes.

Rather than treating these records as
independent textual experiences, the graph
explicitly associates them with a shared
architectural relationship while preserving
their individual source records.

Importantly, the graph contains additional
connections involving executed designs,
architectural parameters, measured metrics,
and Reflection records.

The complete snapshot after Iteration 3
contains 330 typed edges. These represent
connections among graph records, not
330 independent architectural relationships
or simulation observations.

This connected representation enables
V3 to accumulate evidence across iterations
and maintain an explicit assessment of
each recorded parameter--metric relationship.

The following step examines how the
available runtime evidence determines
the relationship's recorded state
and causal-strength assessment.


% =====================================================
% STEP 5: EVIDENCE STATE
% =====================================================

\subsubsection{Step 5: Interpreting Relationship Evidence}

The Runtime Evidence Graph does not treat
every recorded relationship as established
architectural knowledge.

Instead, each relationship maintains
an evidence-based assessment derived
from the available runtime observations.

For the selected L1D-size--area relationship,
the recorded evidence contains one
supporting observation and two records
classified as insufficient evidence.

No exact parameter-isolated comparisons
are available, and the graph records
three confounded observations.

Table~\ref{tab:app_v3_relationship_state}
summarizes the relationship assessment.

\begin{table}[htbp]
\centering
\caption{Recorded evidence assessment for
the L1D-size--area relationship.}
\label{tab:app_v3_relationship_state}
\renewcommand{\arraystretch}{1.15}
\small
\begin{tabular}{lc}
\toprule
\textbf{Evidence Property} &
\textbf{Recorded Value} \\
\midrule
Source iterations & 2, 3 \\
Linked claims & 3 \\
Linked evidence records & 3 \\
Supporting evidence & 1 \\
Insufficient evidence & 2 \\
Exact parameter isolations & 0 \\
Confounded observations & 3 \\
Relationship state &
\texttt{provisionally\_supported} \\
Causal strength &
\texttt{confounded} \\
\bottomrule
\end{tabular}
\end{table}

The relationship is classified as
\texttt{provisionally\_supported}
because one independent runtime iteration
provides supporting evidence.

However, its causal strength remains
\texttt{confounded}, since the available
comparisons involve changes to multiple
architectural parameters.

Consequently, the observed relationship
cannot be interpreted as demonstrating
that reducing L1D cache size alone
causes area to decrease.

This example illustrates an important
property of V3: the graph can record
supporting evidence while simultaneously
preserving uncertainty about the
underlying causal interpretation.


% =====================================================
% STEP 6: COMPARISON WITH DYNAMIC MEMORY
% =====================================================

% =====================================================
% STEP 6: DYNAMIC MEMORY VS RUNTIME EVIDENCE GRAPH
% =====================================================

\subsubsection{Step 6: Comparing Graph Organization with Dynamic Memory}

The principal difference between V2 and V3
concerns how accumulated architectural
evidence is organized and made available
during subsequent candidate generation.

In V2, Dynamic Memory stores previous DSE
experiences in an iteration-based structure.
Each experience contains the executed design,
simulation results, extracted claims,
runtime evidence, Reflection feedback,
and ExpeL lessons.

In V3, these information types are represented
through explicit graph nodes and edges.
Claims describing the same architectural
relationship are associated with a shared
parameter--metric relationship node that
maintains aggregated evidence assessments.

To illustrate this difference, the same
architectural prediction is examined using
recorded V2 and V3 experiments:

\begin{equation}
\mathrm{L1D\ Size}\downarrow
\;\rightarrow\;
\mathrm{Area}\downarrow.
\end{equation}


% =====================================================
% V2: DYNAMIC MEMORY
% =====================================================

\paragraph{V2: Iteration-Based Dynamic Memory.}

In the recorded V2 experiment
\texttt{run\_20260726\_212046},
the same architectural prediction appears
in four separate iterations.

The Claim Extractor repeatedly identifies
the following prediction:

\begin{quote}
\emph{``Reducing l1d\_size may decrease area.''}
\end{quote}

Listing~\ref{lst:app_v2_repeated_claims}
presents selected fields from the actual
Dynamic Memory records.

\begin{lstlisting}[
    style=remodeljson,
    caption={Selected Dynamic Memory claim records
    from Iterations 2--5 in V2.},
    label={lst:app_v2_repeated_claims},
    breaklines=true
]
{
  "experiences": [
    {
      "iteration": 2,
      "extracted_claims": [
        {
          "claim_text": "Reducing l1d_size may decrease area.",
          "claim_type": "verifiable_claim",
          "reasoning_type": "parameter_change",
          "design_parameter_changed": "l1d_size",
          "parameter_direction": "decrease",
          "old_value": null,
          "new_value": "16kB",
          "metric": "area",
          "expected_direction": "decrease",
          "reason_not_verifiable": null,
          "source_section": "analysis.steps[0]"
        }
      ]
    },
    {
      "iteration": 3,
      "extracted_claims": [
        {
          "claim_text": "Reducing l1d_size may decrease area.",
          "new_value": "8kB",
          "metric": "area",
          "expected_direction": "decrease"
        }
      ]
    },
    {
      "iteration": 4,
      "extracted_claims": [
        {
          "claim_text": "Reducing l1d_size may decrease area.",
          "new_value": "16kB",
          "metric": "area",
          "expected_direction": "decrease"
        }
      ]
    },
    {
      "iteration": 5,
      "extracted_claims": [
        {
          "claim_text": "Reducing l1d_size may decrease area.",
          "new_value": "16kB",
          "metric": "area",
          "expected_direction": "decrease"
        }
      ]
    }
  ]
}
\end{lstlisting}

As shown in Listing~\ref{lst:app_v2_repeated_claims},
the same parameter--metric prediction is
preserved within four different
iteration-based experiences.

Although the predicted relationship
remains the same, the proposed parameter
values vary across iterations.

For example, Iteration 2 proposes an
L1D cache size of 16 kB, while
Iteration 3 proposes 8 kB.

Dynamic Memory retains these claims
alongside their corresponding evidence,
Reflection feedback, and experiential
lessons.

However, the repeated predictions
remain organized primarily within
their original iteration records,
rather than being explicitly associated
with one shared evidence-assessed
relationship.


% =====================================================
% V3: RUNTIME EVIDENCE GRAPH
% =====================================================

\paragraph{V3: Relationship-Centered Runtime Evidence Graph.}

In contrast, the recorded V3 experiment
\texttt{run\_20260731\_195619}
organizes repeated architectural
predictions through shared
parameter--metric relationships.

The graph snapshot after Iteration 3
contains the relationship:

\begin{quote}
\texttt{relationship::l1d\_size::decrease::area::decrease}
\end{quote}

Three extracted claims generated during
Iterations 2 and 3 contribute to this
single relationship record.

Listing~\ref{lst:app_v3_shared_relationship}
presents selected fields from the
recorded relationship node.

\begin{lstlisting}[
    style=remodeljson,
    caption={Recorded V3 relationship node
    linking repeated L1D-size--area claims.},
    label={lst:app_v3_shared_relationship},
    breaklines=true
]
{
  "node_id": "relationship::l1d_size::decrease::area::decrease",
  "node_type": "parameter_metric_relationship",
  "parameter": "l1d_size",
  "parameter_direction": "decrease",
  "metric": "area",
  "expected_metric_direction": "decrease",

  "claim_ids": [
    "claim::iteration_2::0",
    "claim::iteration_3::0",
    "claim::iteration_3::11"
  ],

  "evidence_ids": [
    "evidence::iteration_2::0",
    "evidence::iteration_3::0",
    "evidence::iteration_3::11"
  ],

  "source_iterations": [2, 3],

  "evidence_counts": {
    "supported": 1,
    "contradicted": 0,
    "unchanged": 0,
    "mixed": 0,
    "insufficient_evidence": 2,
    "non_verifiable": 0
  },

  "evidence_quality": {
    "exact_parameter_isolations": 0,
    "confounded_observations": 3
  },

  "relationship_state": "provisionally_supported",
  "causal_strength": "confounded"
}
\end{lstlisting}

Unlike the iteration-based organization
in V2, the V3 graph explicitly connects
the three claims to a shared
relationship record.

The \texttt{claim\_ids} field identifies
the individual claims associated with
the relationship, while
\texttt{evidence\_ids} preserves links
to their corresponding runtime evidence.

The \texttt{source\_iterations} field
identifies the iterations from which
the accumulated evidence originates.

The graph also maintains an explicit
evidence assessment.

In this example, one supporting
observation and two insufficient-evidence
records are associated with the
relationship.

The resulting state is
\texttt{provisionally\_supported},
while its causal strength remains
\texttt{confounded} because no exact
parameter-isolated comparisons
are available.

Therefore, the graph does not
automatically treat repeated
architectural predictions as
established causal knowledge.


% =====================================================
% INTERPRETATION
% =====================================================

% =====================================================
% INTERPRETATION
% =====================================================


% =====================================================
% STEP 7: GRAPH-INFORMED ACG REASONING
% =====================================================

\subsubsection{Step 7: Graph-Informed ACG Reasoning}

The preceding steps demonstrated how V3
organizes accumulated architectural claims
and runtime evidence into explicit
parameter--metric relationships.

This step examines how the resulting
relationship information is supplied
to ACG and reflected in subsequent
architectural reasoning.

The example continues with
\texttt{run\_20260731\_195619},
using the graph snapshot after
Iteration 3 and the ACG reasoning
generated during Iteration 4.

\paragraph{Graph Evidence Supplied to ACG.}

Before generating the next architectural
candidate, V3 constructs a compact
graph-informed refinement context.

The context contains the accumulated
parameter--metric relationships,
their evidence assessments, deterministic
rankings, and relationship-specific
evidence details.

For the previously examined relationship,

\begin{equation}
\mathrm{L1D\ Size}\downarrow
\;\rightarrow\;
\mathrm{Area}\downarrow,
\end{equation}

Listing~\ref{lst:app_v3_forwarded_relationship}
presents selected fields from the
actual Iteration 4 ACG input.

\begin{lstlisting}[
    style=remodeljson,
    caption={Recorded compact graph relationship
    forwarded to ACG during Iteration 4.},
    label={lst:app_v3_forwarded_relationship},
    breaklines=true
]
{
  "node_id": "relationship::l1d_size::decrease::area::decrease",
  "parameter": "l1d_size",
  "parameter_direction": "decrease",
  "metric": "area",
  "expected_direction": "decrease",
  "relationship_state": "provisionally_supported",
  "causal_strength": "confounded",
  "source_iterations": [2, 3],
  "evidence_counts": {
    "supported": 1,
    "contradicted": 0,
    "unchanged": 0,
    "mixed": 0,
    "insufficient_evidence": 2,
    "non_verifiable": 0
  },
  "num_independent_runtime_iterations": 1,
  "rank": 1,
  "confidence_score": 0.7175,
  "priority_score": 0.6457,
  "recommended_action": "targeted_validation"
}
\end{lstlisting}

The forwarded record identifies the
relationship as
\texttt{provisionally\_supported},
with \texttt{causal\_strength=confounded}.

Although the relationship is ranked
first among the available relationships,
its recommended action is
\texttt{targeted\_validation},
indicating that additional evidence
would be useful before treating the
predicted effect as established.

The associated evidence-detail record
also reports zero exact parameter
isolations and three confounded
observations.

Furthermore, its reasoning policy states:

\begin{quote}
\emph{``One independent iteration supports
this relationship. Describe it as
provisional rather than established.''}
\end{quote}

Thus, the information supplied to ACG
includes both the predicted architectural
relationship and the limitations of
its supporting evidence.


\paragraph{Recorded ACG Reasoning.}

During Iteration 4, the selected ACG
candidate proposed an architectural
configuration containing an L1D cache
size of 8 kB.

The recorded reasoning included:

\begin{quote}
\emph{``Reducing L1D size can help lower
the area metric, which is currently
within bounds but could be improved
for better cost feasibility.''}
\end{quote}

The corresponding conclusion stated:

\begin{quote}
\emph{``By decreasing l1d\_size, we expect
to decrease the area, helping to
maintain compliance with cost constraints.''}
\end{quote}

The final explanation further stated:

\begin{quote}
\emph{``This proposal is based on mixed
runtime evidence that supports the idea
of reducing the L1D size to improve cost
metrics while maintaining the robustness
of performance metrics.''}
\end{quote}

These statements are consistent with
the graph's recorded L1D-size--area
relationship and indicate that the
generated explanation refers to
previous runtime evidence.

However, the explanation does not
explicitly identify the relationship
as \texttt{provisionally\_supported}
or report its confounded causal strength.

Furthermore, although the mandatory
graph-grounding instructions require
the LLM to report the relationship
identifier, evidence state, causal
strength, and source iterations,
these fields are not explicitly
reported in the selected reasoning.

The generated explanation also refers
to ``mixed runtime evidence,''
whereas the graph's recorded
relationship state is
\texttt{provisionally\_supported}.
Therefore, this wording should be
interpreted as the LLM's description
rather than the graph's formal
evidence classification.


\paragraph{Interpretation.}

This example demonstrates that V3
successfully makes accumulated
relationship-level evidence and
its uncertainty information available
to subsequent ACG reasoning.

The selected ACG response is consistent
with the recorded relationship and
refers to previous runtime evidence
when motivating the architectural
candidate.

However, the response does not fully
satisfy the mandatory graph-grounding
instructions, particularly the
requirement to explicitly report
relationship metadata and evidence
provenance.

Thus, the example illustrates
a distinction between providing
structured evidence and ensuring
its complete application by the LLM.

The Runtime Evidence Graph establishes
a structured basis for evidence-aware
candidate generation, but the LLM
may still omit important evidence
qualifications or fail to follow
all grounding instructions.

This limitation is consistent with
the broader V3 evaluation, which
distinguishes graph relationship
availability from actual graph-grounded
reasoning behavior.


\subsubsection{Interpretation of the Graph Organization Example}

The recorded examples illustrate the principal
organizational difference between textual
Dynamic Memory in V2 and the Runtime Evidence
Graph in V3.

In V2, repeated architectural predictions
are preserved within their corresponding
iteration-based experiences, alongside
runtime evidence, Reflection feedback,
and ExpeL lessons.

In contrast, V3 retains these individual
records while explicitly connecting
related predictions to shared
parameter--metric relationships.

These relationships maintain accumulated
evidence counts, source iterations,
relationship states, and causal-strength
information, allowing the framework to
distinguish supported observations from
contradictory, confounded, or insufficient
evidence.

Importantly, V3 does not eliminate
repeated claims or their original evidence.
Instead, it consolidates their relationship-level
representation while preserving the
underlying execution records and uncertainty.

Furthermore, the complete Runtime Evidence
Graph remains available internally, while
subsequent ACG stages receive a compact
relationship-level context containing
evidence assessments and evidence-aware
rankings rather than the full textual
experience history.

The V2 and V3 examples originate from
different experimental runs and are
intended to illustrate differences in
evidence organization rather than
differences in architectural performance
or reasoning correctness.

Overall, this case study demonstrates
the structural contribution of V3:
transforming accumulated execution-grounded
information into explicit, reusable,
and evidence-assessed architectural
relationships.

However, this organizational improvement
does not by itself establish that
graph-informed reasoning is more correct
than textual-memory-based reasoning.


The subsequent ACG example further demonstrates
that structured relationship evidence can
inform architectural reasoning, while also
revealing that the LLM does not always
explicitly apply every required grounding
instruction.


